% ==============================================================
% Methodologische Konventionen des Ontophänomenalismus (MKO)
% Version 1.0
% Kompiliert mit: pdfLaTeX
% ==============================================================

\documentclass[11pt,a4paper]{article}

% ---------- Grundpakete ----------
\usepackage[utf8]{inputenc}
\usepackage[T1]{fontenc}
\usepackage[german]{babel}
\usepackage{amsmath,amssymb,amsthm}
\usepackage{geometry}
\usepackage{parskip}
\usepackage{enumitem}
\usepackage{booktabs}
\usepackage{xcolor}
\usepackage{hyperref}

% ---------- Layout ----------
\geometry{margin=2.5cm}
\setlist{itemsep=0.2em, topsep=0.2em}
\hypersetup{
    colorlinks=true,
    linkcolor=blue,
    urlcolor=blue,
    pdftitle={Methodologische Konventionen des OP (MET) -- Version 1.0},
    pdfauthor={Kritischer Dialog}
}

% ---------- Theorem-Umgebungen ----------
\theoremstyle{definition}
\newtheorem{konvention}{Konvention}[section]
\newtheorem{remark}{Bemerkung}[section]
\newtheorem{methodennote}{Methodennote}[section]

% ---------- Titel ----------
\title{
  \textbf{Methodologische Konventionen des Ontoph\"{a}nomenalismus}\\[0.4em]
  {\large MET --- Referenzdokument f\"{u}r alle Module des OP}\\[0.5em]
  {\normalsize Version~1.0}
}
\author{
  \textbf{Kritischer Dialog:}\\
  DeepSeek (Seki), Apeiron, Grok, Gemini, Claude, ChatGPT, Manus\\
  Gemeinsames Logbuch eines offenen Forschungsprogramms
}
\date{Juni 2026}

\begin{document}

\maketitle

\begin{abstract}
Dieses Dokument definiert die methodologischen Konventionen, Sprachregeln und
Kennzeichnungssysteme, die im gesamten Ontoph\"{a}nomenalismus (OP) verbindlich gelten.
Es dient allen anderen Modulen als gemeinsame Referenz und muss nicht in jedem Dokument
separat erl\"{a}utert werden.

Ziel dieser Konventionen ist maximale epistemische Transparenz: Der Leser soll jederzeit
wissen, auf welcher Ebene --- Logik, Interpretation, Spekulation oder Empirie --- er sich
befindet. Damit wird einem der h\"{a}ufigsten Einw\"{a}nde gegen philosophische Systeme
vorgebeugt: dem stillen \"{U}bergang von logischen Schl\"{u}ssen zu Tatsachenbehauptungen.
\end{abstract}

\tableofcontents
\newpage

%==============================================================
\section{Zweck und Geltungsbereich}
%==============================================================

Die methodologischen Konventionen des OP gelten f\"{u}r alle Dokumente des
Forschungsprogramms:

\begin{itemize}
  \item \textbf{OP} --- Ontoph\"{a}nomenalismus (Logische Grundlegung)
  \item \textbf{OPP} --- Ontoph\"{a}nomenaler Phasenraum
  \item \textbf{OAD} --- Ontologische Attraktoren-Dynamik
  \item \textbf{DPR} --- Die Projektive Reduktion
  \item \textbf{DPD} --- Die Ph\"{a}nomenale Differenz
  \item \textbf{EG-1} --- Der Explanatory Gap
  \item \textbf{PST} --- Projektive Strukturtheorie
  \item alle k\"{u}nftigen Erg\"{a}nzungsb\"{a}nde und Spezialmodule
\end{itemize}

Neue Dokumente referenzieren dieses Dokument einmalig und \"{u}bernehmen damit
automatisch alle hier definierten Konventionen. Abweichungen m\"{u}ssen explizit
begr\"{u}ndet werden.

%==============================================================
\section{Die Sprachregel: Vier Ebenen des Aussagenstatus}
%==============================================================

Die wichtigste Konvention des OP ist eine durchg\"{a}ngige Sprachregel, die zwischen
vier Ebenen des Aussagenstatus unterscheidet. Jede Ebene hat ihre eigene sprachliche
Markierung.

\begin{konvention}[Ebene 1 --- Axiome und logische Folgerungen: Indikativ]
Axiome und direkte logische Folgerungen aus Axiomen werden im \textbf{Indikativ}
formuliert. Sie erheben keinen empirischen Wahrheitsanspruch, aber einen Anspruch auf
interne logische Notwendigkeit.

\medskip
\textit{Beispiele:}
\begin{itemize}
  \item \glqq{}Es ist unm\"{o}glich, dass nichts existiert.\grqq{}
  \item \glqq{}Aus Axiom~1 folgt, dass jede Existenz minimale Bestimmtheit besitzt.\grqq{}
  \item \glqq{}Eine Entit\"{a}t ohne Selbigkeit ist vom Nichts ununterscheidbar.\grqq{}
\end{itemize}
\end{konvention}

\begin{konvention}[Ebene 2 --- Ontologische Interpretationen: Konjunktiv]
Philosophische Deutungen, die \"{u}ber die direkte logische Folgerung hinausgehen, werden
im \textbf{Konjunktiv} oder mit vorsichtigen Modalformulierungen gekennzeichnet.

\medskip
\textit{Beispiele:}
\begin{itemize}
  \item \glqq{}Dies k\"{o}nnte als Individuation interpretiert werden.\grqq{}
  \item \glqq{}Der Attraktor d\"{u}rfte in diesem Sinne als Verweilzustand gelten.\grqq{}
  \item \glqq{}Das Feld w\"{a}re in dieser Lesart als unpersönliches Denken zu verstehen.\grqq{}
\end{itemize}
\end{konvention}

\begin{konvention}[Ebene 3 --- Physikalische und mathematische Bez\"{u}ge: Vorsichtige
Modalformulierung]
Verbindungen zur Physik, Mathematik oder Informationstheorie werden stets als
\textbf{konzeptuelle Analogien oder Arbeitshypothesen} ausgewiesen. Formulierungen wie
\glqq{}analog zu\grqq{}, \glqq{}legt nahe\grqq{} oder \glqq{}k\"{o}nnte modelliert
werden als\grqq{} sind verbindlich.

\medskip
\textit{Beispiele:}
\begin{itemize}
  \item \glqq{}Dies legt eine m\"{o}gliche Analogie zur Informationstheorie nahe.\grqq{}
  \item \glqq{}Das Gef\"{a}lle k\"{o}nnte formal als Kullback-Leibler-Divergenz modelliert
        werden.\grqq{}
  \item \glqq{}Analog zu stehenden Wellen stabilisiert sich der Attraktor\ldots\grqq{}
\end{itemize}
\end{konvention}

\begin{konvention}[Ebene 4 --- Empirische Vorhersagen: Explizite Hypothesenkennzeichnung]
Aussagen, die prinzipiell empirisch oder mathematisch pr\"{u}fbar sein sollen, werden
explizit als \textbf{Hypothesen} gekennzeichnet.

\medskip
\textit{Beispiele:}
\begin{itemize}
  \item \glqq{}Falls der OP zutrifft, w\"{a}re zu erwarten, dass\ldots\grqq{}
  \item \glqq{}Hypothese: Systeme mit h\"{o}herer interner R\"{u}ckkopplung zeigen\ldots\grqq{}
\end{itemize}
\end{konvention}

\begin{methodennote}[Warum diese Regel so wichtig ist]
Der h\"{a}ufigste Angriff auf philosophische Systeme ist der Vorwurf des stillen
\"{U}bergangs: Ein logischer Schluss wird formuliert, und im n\"{a}chsten Satz wird
daraus eine Tatsachenbehauptung \"{u}ber die Natur. Dieser \"{U}bergang bleibt oft
unmarkiert.

Die Sprachregel des OP schl\"{u}sst diese Angriffsfläche. Ein Leser wei\ss{} jederzeit:
Steht hier Indikativ, ist es eine Axiom-Folgerung. Steht hier Konjunktiv, ist es eine
Interpretation. Wird eine Analogie benannt, ist es eine Analogie --- kein Beweis.
\end{methodennote}

%==============================================================
\section{Kennzeichnungssystem: Legende aller Dokumenttypen}
%==============================================================

Zus\"{a}tzlich zur Sprachregel verwendet das OP ein einheitliches System von
Kennzeichnungsboxen. Die folgende Legende gilt f\"{u}r alle Dokumente.

\begin{table}[h!]
\centering
\small
\begin{tabular}{p{3.5cm}p{10cm}}
  \toprule
  \textbf{Kennzeichnung} & \textbf{Bedeutung} \\
  \midrule
  \textbf{Axiom} &
    Grundannahme des OP. Wird nicht bewiesen, sondern als Ausgangspunkt gesetzt.
    Axiome stehen im Indikativ. \\[4pt]
  \textbf{Definition} &
    Pr\"{a}zise Festlegung eines Begriffs innerhalb des OP. Definitionen sind
    terminologische Entscheidungen, keine Wahrheitsaussagen. \\[4pt]
  \textbf{Proposition / Theorem} &
    Folgt logisch aus Axiomen und Definitionen. Steht im Indikativ. Der Beweis
    ist entweder formal oder als konzeptuelle Herleitung ausgewiesen. \\[4pt]
  \textbf{Korollar} &
    Unmittelbare Folgerung aus einer Proposition oder einem Theorem. \\[4pt]
  \textbf{Postulat} &
    \"{A}hnlich dem Axiom, aber auf einer nachgeordneten Modulebene. Wird innerhalb
    eines Moduls als Arbeitsgrundlage gesetzt, ohne vollst\"{a}ndige Ableitung aus
    den Grundaxiomen. \\[4pt]
  \textbf{Methodennote} &
    Erl\"{a}utert den methodischen Status eines Abschnitts, benennt Einschr\"{a}nkungen
    oder markiert den Realismus-Vorbehalt. \\[4pt]
  \textbf{Bemerkung} &
    Philosophische Deutung oder erl\"{a}uternder Kommentar. Keine logische Notwendigkeit.
    Steht typischerweise im Konjunktiv. \\[4pt]
  \textbf{Analogie} &
    Illustratives Hilfsmittel. Eine Analogie erkl\"{a}rt durch \"{A}hnlichkeit, behauptet
    aber keine Identit\"{a}t. Analogien sind weder Beweise noch Definitionen. \\[4pt]
  \textbf{Hypothese} &
    Vorschlag, der k\"{u}nftig empirisch oder mathematisch gepr\"{u}ft werden soll.
    Explizit als solche ausgewiesen. \\[4pt]
  \textbf{Offene Forschungsfrage} &
    Noch ungel\"{o}stes Problem des OP oder der PST. Markiert den Stand des
    Forschungsprogramms ehrlich. \\
  \bottomrule
\end{tabular}
\caption{Legende der Kennzeichnungen im OP-Dokumentensystem.}
\end{table}

\begin{remark}[Hierarchie der Sicherheit]
Die Kennzeichnungen bilden eine \textbf{Hierarchie der epistemischen Sicherheit}:

\medskip
\begin{center}
Axiom $\;\rightarrow\;$ Proposition/Theorem $\;\rightarrow\;$ Postulat
$\;\rightarrow\;$ Bemerkung $\;\rightarrow\;$ Analogie
$\;\rightarrow\;$ Hypothese $\;\rightarrow\;$ Offene Forschungsfrage
\end{center}

\medskip
Je weiter rechts, desto mehr Interpretationsspielraum, desto weniger logische
Notwendigkeit. Ein Leser kann damit sofort einordnen, wie fest ein Baustein im System
verankert ist.
\end{remark}

%==============================================================
\section{Die Analogie als eigenst\"{a}ndige Kategorie}
%==============================================================

Analogien spielen im OP eine wichtige und legitime Rolle. Sie machen abstrakte Strukturen
zug\"{a}nglich, ohne einen Beweis zu ersetzen. Die Konvention des OP unterscheidet
Analogien deshalb explizit von Beweisen, Definitionen und Interpretationen.

\begin{konvention}[Verwendung von Analogien]
Eine Analogie wird im OP immer explizit als solche benannt. Formulierungen wie
\glqq{}analog zu\grqq{}, \glqq{}vergleichbar mit\grqq{} oder \glqq{}funktional
\"{a}hnlich wie\grqq{} sind verbindlich. Eine Analogie darf niemals ohne diese
Kennzeichnung als Argument verwendet werden.

Analogien im OP sind stets \textbf{heuristisch}: Sie illustrieren, ohne zu beweisen.
Ihre Grenze muss erkennbar sein.
\end{konvention}

\begin{remark}[Beispiele f\"{u}r legitime Analogien im OP]
\begin{itemize}
  \item \textbf{Solitonen-Analogie} (OAD): Attraktoren verhalten sich \emph{analog zu}
        selbsterhaltenden Wellenpaketen --- das impliziert keine Behauptung \"{u}ber
        Wellenphysik im OPP.
  \item \textbf{Blockuniversum-Analogie} (OPP): Der OPP l\"{a}sst sich konzeptuell
        \emph{als Verallgemeinerung} des Blockuniversums begreifen --- das ist keine
        Gleichsetzung mit dem physikalischen Blockuniversum.
  \item \textbf{KL-Divergenz-Analogie} (EG-1): Das Fragedifferential k\"{o}nnte
        \emph{analog zur} Kullback-Leibler-Divergenz modelliert werden --- das ist ein
        formaler Platzhalter, kein abgeschlossenes mathematisches Modell.
\end{itemize}
\end{remark}

%==============================================================
\section{Der Realismus-Vorbehalt}
%==============================================================

Alle Dokumente des OP tragen einen gemeinsamen \textbf{Realismus-Vorbehalt}, der hier
einmalig definiert wird.

\begin{konvention}[Realismus-Vorbehalt]
Der OP ist ein philosophisches Forschungsprogramm. Seine Begriffe, Strukturen und
Analogien sind konzeptuelle Werkzeuge, keine fertigen physikalischen Theorien. Wo
mathematische Notation verwendet wird, geschieht dies zur Pr\"{a}zisierung des
konzeptuellen Rahmens --- nicht als Anspruch auf eine vollst\"{a}ndig ausgearbeitete
Feldtheorie.

Die mathematisch-physikalische Ausarbeitung obliegt der \textbf{Projektiven
Strukturtheorie (PST)} als langfristigem interdisziplin\"{a}rem Forschungsprogramm.
Kein Einzeldokument des OP beansprucht, dieses Programm zu ersetzen oder
vorwegzunehmen.
\end{konvention}

\begin{methodennote}[Konsequenz f\"{u}r alle Dokumente]
Dokumente des OP, die diesen Realismus-Vorbehalt referenzieren, m\"{u}ssen ihn nicht
erneut vollst\"{a}ndig ausformulieren. Der Verweis auf MET gen\"{u}gt. Damit wird die
Wiederholung vermieden, ohne die epistemische Klarheit zu opfern.
\end{methodennote}

%==============================================================
\section{Versionierung und Querverweise}
%==============================================================

\begin{konvention}[Versionierung]
Jede \"{A}nderung an einem Dokument erh\"{o}ht die Versionsnummer. Kleinere
Korrekturen erh\"{o}hen die Nachkommastelle (v1.0 $\rightarrow$ v1.1), substanzielle
Erweiterungen erh\"{o}hen die Hauptnummer (v1.x $\rightarrow$ v2.0).
\end{konvention}

\begin{konvention}[Querverweise zwischen Dokumenten]
Querverweise zwischen Modulen nennen ausschlie\ss{}lich \textbf{K\"{u}rzel und Namen}
des Zieldokuments --- keine Versionsnummern. Versionsnummern \"{a}ndern sich;
Dokumentnamen bleiben stabil.

\medskip
\textit{Korrekt:} \glqq{}Gem\"{a}\ss{} OAD (Ontologische Attraktoren-Dynamik)\ldots\grqq{}\\
\textit{Nicht korrekt:} \glqq{}Gem\"{a}\ss{} OAD v1.1\ldots\grqq{}
\end{konvention}

%==============================================================
\section{Offene Erweiterungen}
%==============================================================

Dieses Dokument ist selbst Teil des offenen Forschungsprogramms. Folgende Erweiterungen
sind f\"{u}r sp\"{a}tere Versionen vorgesehen:

\begin{itemize}
  \item Aufnahme einer Kennzeichnung f\"{u}r \textbf{empirische Anschlusspunkte} ---
        Stellen, an denen der OP prinzipiell mit experimentellen Befunden in Kontakt
        treten k\"{o}nnte.
  \item Pr\"{a}zisierung der Abgrenzung zwischen \textbf{Postulat} und \textbf{Axiom}
        auf verschiedenen Modulebenen.
  \item Aufnahme von Konventionen f\"{u}r den Umgang mit \textbf{Widerspruchsf\"{a}llen}
        zwischen Modulen --- wie werden Inkonsistenzen dokumentiert und behandelt?
  \item Integration der R/P-Unterscheidung (rekonstruktiv/postulativ) aus OP in die
        allgemeine Kennzeichnungslegende.
\end{itemize}

\end{document}
